\chapter{Android приложение и его функционал}
\label{ch:android}

\section{Лабораторная работа №6. Местоположение смартфона}
\subsection{Цель}
Получить доступ к данным о местоположении Android-телефона и вывести на экран значения.

\subsection{Теория}
Для определения местоположения в Android используются GPS, Wi-Fi и сотовые вышки. Класс FusedLocationProviderClient объединяет данные от всех источников, обеспечивая оптимальный баланс точности и энергопотребления. Для доступа к координатам требуется запросить разрешения ACCESS\_FINE\_LOCATION или ACCESS\_COARSE\_LOCATION \cite{access_location}.

\subsection{Реализация}
В проекте создано Activity location.kt. При запуске запрашиваются разрешения, инициализируется FusedLocationProviderClient \cite{fused_location}:

\begin{lstlisting}[language=Java, caption=Инициализация FusedLocationProviderClient]
class location:AppCompatActivity(){
    lateinit var myFusedLocationProviderClient:FusedLocationProviderClient
    
    override fun onCreate(savedInstanceState: Bundle?){
        super.onCreate(savedInstanceState)
        myFusedLocationProviderClient = LocationServices.getFusedLocationProviderClient(this)}}
\end{lstlisting}

Метод getCurrentLocation() получает координаты с высокой точностью:

\begin{lstlisting}[language=Java, caption=Получение текущего местоположения]
private fun getCurrentLocation(){
    if (checkPermissions()){
        if (isLocationEnabled()){
            myFusedLocationProviderClient.getCurrentLocation(Priority.PRIORITY_HIGH_ACCURACY,null).addOnCompleteListener(this){task->
                val location:Location?=task.result
                if (location!=null){
                    lat_site=location.latitude
                    lon_site=location.longitude
                    lat.setText(location.latitude.toString())
                    lon.setText(location.longitude.toString())
                    alt.setText(round(location.altitude).toString())
                    tojson(location.latitude,location.longitude,location.altitude,location.time.toString())}}}}}
\end{lstlisting}

Данные выводятся в TextView, а также сохраняются в файл location.json в папке Downloads. Кнопка card\_button открывает карту в браузере:

\begin{lstlisting}[language=Java, caption=Открытие карты]
card_button.setOnClickListener({
    val uri=Uri.parse("https://geotree.ru/coordinates?lat=$lat_site&lon=$lon_site")
    val siteIntent=Intent(Intent.ACTION_VIEW,uri)
    startActivity(siteIntent)})
\end{lstlisting}

Обновление координат происходит каждые 10 секунд:

\begin{lstlisting}[language=Java, caption=Периодическое обновление]
private fun updatelocation(){
    handler.postDelayed({
        getCurrentLocation()
        updatelocation()}, 10000)}
\end{lstlisting}

\begin{figure}[h]
\centering
\includegraphics[width=0.3\textwidth]{lab6_android.png}
\caption{Экран Location}
\end{figure}

\section{Лабораторная работа №8. Работа с сокетами (ZMQ). Передача данных от Android к PC}
\subsection{Цель}
Реализовать обмен сообщениями между Android-приложением и сервером на ПК с использованием ZeroMQ \cite{zeromq}.

\subsection{Теория}
ZeroMQ - высокопроизводительная асинхронная библиотека для обмена сообщениями, поддерживающая различные паттерны. В отличие от обычных сокетов, ZMQ автоматически управляет буферизацией, переподключением и маршрутизацией. Для взаимодействия Android-клиента с ПК-сервером используется сокет типа REQ (запрос) и REP (ответ).

\subsection{Реализация}
Создано Activity sockets.kt. При запуске запускается поток сервера, который связывается с портом 6666:

\begin{lstlisting}[language=Java, caption=Серверная часть на ZMQ]
fun startServer(){
    val context=ZMQ.context(1)
    val socket=ZContext().createSocket(SocketType.REP)
    socket.bind("tcp://localhost:6666")
    var counter:Int=0
    
    while (true){
        counter++
        val requestBytes=socket.recv(0)
        val request=String(requestBytes,ZMQ.CHARSET)
        handler.postDelayed({serverView.text=request},0)
        val response="Получено сообщение №$counter"
        socket.send(response.toByteArray(ZMQ.CHARSET),0)}}
\end{lstlisting}

Клиент отправляет текст из EditText при нажатии кнопки. Переключатель позволяет менять адрес сервера на внешний IP:

\begin{lstlisting}[language=Java, caption=Клиентская часть]
fun startClient(){
    val context=ZMQ.context(1)
    val socket=ZContext().createSocket(SocketType.REQ)
    if (flag){socket.connect("tcp://192.168.1.211:2222")
    }else{socket.connect("tcp://localhost:6666")}
    val request=clientText.text.toString()
    socket.send(request.toByteArray(ZMQ.CHARSET),0)
    val reply=String(socket.recv(0),ZMQ.CHARSET)
    clientView.text=reply}
\end{lstlisting}

\begin{figure}[h]
\centering
\includegraphics[width=0.3\textwidth]{lab8_android.png}
\caption{Экран Sockets}
\end{figure}

\section{Лабораторная работа №10. Android Cellinfo}
\subsection{Цель}
Получать и передавать на сервер детальную информацию о сотах сотовой связи вместе с геолокацией и сетевым трафиком.

\subsection{Теория}
Класс TelephonyManager \cite{telephony} предоставляет метод getAllCellInfo(), возвращающий список сотовых вышек текущего типа сети (LTE, GSM, NR). Для каждой соты можно получить уникальные идентификаторы (CID, TAC, PCI, EARFCN) и параметры сигнала (RSRP, RSRQ, RSSNR). Эти данные вместе с GPS-координатами позволяют анализировать покрытие сети и качество связи.

\subsection{Реализация}
Создано Activity telephony.kt. Пользователь выбирает тип сети (LTE, GSM, NR) через кнопки:

\begin{lstlisting}[language=Java, caption=Выбор типа сети]
bttn_lte.setOnClickListener({
    flag=1
    CellIdentity.setText("CellIdentityLTE")
    CellSignalStrength.setText("CellSignalStrengthLTE")
    bttn_lte.setBackgroundColor(Color.parseColor("#2d2d2f"))
    bttn_gsm.setBackgroundColor(Color.parseColor("#16c603"))
    bttn_nr.setBackgroundColor(Color.parseColor("#16c603"))})
\end{lstlisting}

Получение информации о сотах LTE:

\begin{lstlisting}[language=Java, caption=Получение данных о сотах LTE]
if (flag==1){
    val telephonyManager=getSystemService(Context.TELEPHONY_SERVICE) as TelephonyManager
    val cellInfoList=telephonyManager.allCellInfo
    cellInfoList?.forEach{cellInfo->
        if (cellInfo is CellInfoLte){
            cellIdentityList.add("Band: ${cellInfo.cellIdentity.bandwidth}")
            cellIdentityList.add("CellIdentity: ${cellInfo.cellIdentity.ci}")
            cellIdentityList.add("EARFCN: ${cellInfo.cellIdentity.earfcn}")
            cellIdentityList.add("PCI: ${cellInfo.cellIdentity.pci}")
            cellIdentityList.add("TAC: ${cellInfo.cellIdentity.tac}")
            cellSignalStrengthList.add("RSRP: ${cellInfo.cellSignalStrength.rsrp}")
            cellSignalStrengthList.add("RSRQ: ${cellInfo.cellSignalStrength.rsrq}")
            cellSignalStrengthList.add("RSSI: ${cellInfo.cellSignalStrength.rssi}")}}}
\end{lstlisting}

Данные отображаются в двух списках (идентификаторы и мощность):

\begin{lstlisting}[language=Java, caption=Отображение данных в ListView]
val a1=ArrayAdapter(this,android.R.layout.simple_list_item_1,cellIdentityList)
list_cell_info1.adapter=a1
val a2=ArrayAdapter(this,android.R.layout.simple_list_item_1,cellSignalStrengthList)
list_cell_info2.adapter=a2
\end{lstlisting}

\begin{figure}[h]
\centering
\includegraphics[width=0.3\textwidth]{lab10_android.png}
\caption{Экран Cellinfo}
\end{figure}

\section{Лабораторная работа №11. Android background service}
\subsection{Цель}
Реализовать фоновый сервис, который в фоне собирает данные о местоположении, сотах и трафике, отправляет их на сервер и сохраняет в файл.

\subsection{Теория}
Сервисы Android позволяют приложению выполнять код в фоне, даже когда экран выключен или пользователь работает с другим приложением. Это нужно для периодического сбора местоположения, данных о сотах и отправки их на сервер без участия пользователя.

\subsection{Реализация}
Класс BackgroundService.kt наследуется от Service. В методе onStartCommand запускается корутина, в которой циклически вызываются getCurrentLocation() и getCellInfo():

\begin{lstlisting}[language=Java, caption=Фоновый сервис - основной цикл]
class BackgroundService:Service(){
    private val serviceJob=Job()
    private val serviceScope=CoroutineScope(Dispatchers.IO+serviceJob)
    
    override fun onStartCommand(intent:Intent?,flags:Int,startId:Int):Int{
        flag=intent?.getIntExtra("flag", 1)?:1
        serviceScope.launch{
            while (isActive){
                getCurrentLocation()
                getCellInfo()
                delay(3000)}}
        return START_STICKY}}
\end{lstlisting}

Отправка данных в Activity через LocalBroadcastManager:

\begin{lstlisting}[language=Java, caption=Отправка данных в Activity]
private fun sendLocationToActivity(location:Location){
    val intent=Intent("LOCATION_UPDATE")
    intent.putExtra("latitude",location.latitude.toString())
    intent.putExtra("longitude",location.longitude.toString())
    intent.putExtra("altitude",round(location.altitude).toString())
    intent.putExtra("accuracy",location.accuracy.toString())
    val totalTraffic=TrafficStats.getTotalRxBytes()+TrafficStats.getTotalTxBytes()
    intent.putExtra("traffic","${totalTraffic/1024/1024} MB")
    LocalBroadcastManager.getInstance(this).sendBroadcast(intent)}
\end{lstlisting}

Формирование JSON и отправка на сервер через ZMQ:

\begin{lstlisting}[language=Java, caption=Отправка данных на сервер]
private fun sendAllData(location:Location){
    val jsonData=JSONObject()
    val locData=JSONObject()
    locData.put("latitude",location.latitude)
    locData.put("longitude",location.longitude)
    locData.put("altitude",round(location.altitude))
    jsonData.put("location",locData)
    val cellsArray=JSONArray()
    jsonData.put("cells",cellsArray)
    Thread{
        val socket=ZContext().createSocket(SocketType.REQ)
        socket.connect("tcp://172.22.237.35:443")
        socket.send(jsonData.toString().toByteArray(ZMQ.CHARSET),0)
        socket.close()}.start()
    saveToFile(jsonData)}
\end{lstlisting}

Сохранение данных в файл:

\begin{lstlisting}[language=Java, caption=Сохранение в JSON файл]
private fun saveToFile(jsonData:JSONObject){
    val downloadsDir=Environment.getExternalStoragePublicDirectory(Environment.DIRECTORY_DOWNLOADS)
    val file=File(downloadsDir,"mydata.json")
    val jsonObject=if (file.exists()) JSONObject(file.readText()) else JSONObject()
    val dataArray=if (jsonObject.has("data")) jsonObject.getJSONArray("data") else JSONArray()
    dataArray.put(jsonData)
    jsonObject.put("data",dataArray)
    file.writeText(jsonObject.toString())}
\end{lstlisting}

Запуск и остановка сервиса из Activity:

\begin{lstlisting}[language=Java, caption=Управление сервисом]
bStart.setOnClickListener{
    val startBgIntent=Intent(this,BackgroundService::class.java)
    startBgIntent.putExtra("flag",flag)
    startService(startBgIntent)}

bStop.setOnClickListener{
    val stopBgIntent=Intent(this,BackgroundService::class.java)
    stopService(stopBgIntent)}
\end{lstlisting}

\begin{figure}[h]
\centering
\includegraphics[width=0.3\textwidth]{lab11_android.png}
\caption{Экран location+cellinfo+traffic}
\end{figure}

\endinput